\section{Reflection, Clausen formulas, half shifts, and absolute values}\label{sec:reflection}

\subsection{Kinkelin's reflection formula}

Let
\[
 \Phi(z)=\Log G(1-z)-\Log G(1+z)+z\log(2\pi)
          -\int_0^z\pi x\cot(\pi x)\dd x.
\]
On a simply connected domain avoiding the integers and the zeros of the two
Barnes factors, differentiate and use \cref{eq:phi-global}:
\begin{align*}
 \Phi'(z)
 &=-\phi(1-z)-\phi(1+z)+\log(2\pi)-\pi z\cot(\pi z)\\
 &=z\bigl(\psi(1-z)-\psi(1+z)\bigr)+1-\pi z\cot(\pi z).
\end{align*}
The digamma recurrence and reflection formulas are
\[
 \psi(1+z)=\psi(z)+\frac1z,
 \qquad
 \psi(1-z)-\psi(z)=\pi\cot(\pi z).
\]
They give \(\Phi'(z)=0\).  Since \(\Phi(0)=0\), we obtain:

\begin{theorem}[Kinkelin reflection]\label{thm:kinkelin}
With logarithms continued coherently from \(z=0\),
\begin{equation}\label{eq:kinkelin}
 \Log G(1-z)
 =\Log G(1+z)-z\log(2\pi)
  +\int_0^z\pi x\cot(\pi x)\dd x.
\end{equation}
On a certified real interval \(-1<z<1\) this uses the ordinary real branch
away from zero crossings.  On a complex domain it is a path-dependent
logarithmic identity whose exponential is branch independent.
\end{theorem}

\begin{proof}[Constructive content]
The Gamma reflection and recurrence identities already have finite sine,
product and apartness contracts in the Gamma development.  Cover the input path
by rational disks separated from the integers.  On each disk, evaluate the
derivative of \(\Phi\) in cross-multiplied form, so no unguarded cotangent or
logarithmic derivative is used.  The resulting center is exactly zero; the
radius is made smaller than the local demand.  Integrating the zero derivative
along finitely many segments and using \(\Phi(0)=0\) gives the requested
precision.  Alternatively, expand both sides into their absolutely convergent
Taylor series on \(|z|<1\) and then continue by overlapping disks.  Every step
has a finite branch and separation certificate.
\end{proof}

\subsection{Clausen forms}

For \(0<\theta<2\pi\), define the Clausen function by either of the equivalent
constructive names
\begin{equation}\label{eq:clausen}
 \Cl(\theta)=\sum_{n=1}^{\infty}\frac{\sin(n\theta)}{n^2}
 =-\int_0^\theta\log\left(2\sin\frac t2\right)\dd t.
\end{equation}
The sine series has the uniform tail
\[
 \left|\sum_{n>N}\frac{\sin(n\theta)}{n^2}\right|\le\frac1N.
\]
On a rational compact subinterval of \((0,2\pi)\), integration of the finite
Fourier sum and the same tail proves the equality of the two names.

Integration by parts gives, for \(0<z<1\),
\begin{align}
 \int_0^z\pi x\cot(\pi x)\dd x
 &=z\log(2\sin\pi z)
   +\frac{1}{2\pi}\Cl(2\pi z).
 \label{eq:cot-clausen}
\end{align}
Indeed, the endpoint \(x\log\sin(\pi x)\) tends to zero at \(0\); this limit has
the explicit bound
\(|x\log x|\le C_p2^{-p}\) once \(x\le2^{-M(p)}\), obtained from the log
series and dyadic search.

Substitute \cref{eq:cot-clausen} into \cref{eq:kinkelin}:
\begin{equation}\label{eq:clausen-reflection}
 2\pi\log\frac{G(1-z)}{G(1+z)}
 =2\pi z\log\left(\frac{\sin\pi z}{\pi}\right)
  +\Cl(2\pi z),
 \qquad 0<z<1.
\end{equation}
Using \(G(1+z)=\Gamma(z)G(z)\), this is equivalent to
\begin{equation}\label{eq:clausen-reflection2}
 \log\frac{G(1-z)}{G(z)}
 =z\log\left(\frac{\sin\pi z}{\pi}\right)
  +\log\Gamma(z)+\frac{1}{2\pi}\Cl(2\pi z).
\end{equation}
For iPRA, the primary form is exponentiated and cross-multiplied.  The real
interval hypothesis supplies positive lower bounds for \(G(z),G(1-z)\),
\(\Gamma(z)\), and \(\sin(\pi z)\) once the input box is separated from the
endpoints.

\subsection{The half-shift tangent identity}

Put \(u=\tfrac12-z\) in \cref{eq:kinkelin}, use the Gamma reflection formula,
and simplify.  On \(|z|<\tfrac12\), where \(\cos(\pi z)>0\), one convenient
form is
\begin{equation}\label{eq:halfshift}
 \frac{G(\frac12+z)}{G(\frac12-z)}
 =\frac{(2\pi)^z}{\Gamma(\frac12+z)}
  \sqrt{\frac{\pi}{\cos(\pi z)}}
  \exp\!\left(\pi\int_0^z x\tan(\pi x)\dd x\right).
\end{equation}
An equivalent Bernoulli-polynomial form is
\begin{align}
 \log\frac{G(\frac12+z)}{G(\frac12-z)}
 ={}&\log\Gamma(\tfrac12-z)
 +(z-\tfrac12)\log(2\pi)+\tfrac12\log2\notag\\
 &+\pi\int_0^z(x-\tfrac12)\tan(\pi x)\dd x.
 \label{eq:halfshift-B1}
\end{align}
To verify equivalence, integrate
\(\frac{d}{dx}\log\cos(\pi x)=-\pi\tan(\pi x)\) and use
\[
 \Gamma(\tfrac12+z)\Gamma(\tfrac12-z)=\frac{\pi}{\cos(\pi z)}.
\]

The formal domain is not merely decorative.  A square-root evaluator needs a
positive lower bound for \(\cos(\pi z)\) on real inputs, or a branch path for
complex inputs.  On a rational interval
\(|z|\le\tfrac12-\delta\), elementary sine/cosine bounds give a primitive-
recursive positive lower bound depending on \(\delta\).  The tangent integral
is then ordinary certified quadrature.

\subsection{Conjugation and the squared absolute-value product}\label{sec:modulus}

By \cref{eq:conj},
\begin{equation}\label{eq:modsquare-basic}
 |G(z)|^2=G(z)G(\overline z).
\end{equation}
This identity is preferable to an unsquared modulus in constructive arithmetic:
it avoids square roots and sign choices.

For real \(x\notin\{0,-1,-2,\ldots\}\), divide the canonical products at
\(x+\ii y\), \(x-\ii y\), and \(x\), and simplify each conjugate pair.  The
finite algebra yields
\begin{align}
 |G(x+\ii y)|^2
 ={}&G(x)^2\exp((1+\gamma)y^2)
 \left(1+\frac{y^2}{x^2}\right)\notag\\
 &\times\prod_{k=1}^{\infty}
 \left(1+\frac{y^2}{(x+k)^2}\right)^{k+1}
 \exp\!\left(-\frac{y^2}{k}\right).
 \label{eq:modsquare}
\end{align}
The logarithm of the tail again begins at order \(k^{-2}\).  On a box
\(|x|,|y|\le R\) separated by \(\delta>0\) from the forbidden integers, expand
\(\log(1+y^2/(x+k)^2)\); after subtracting \(y^2/k\), the apparent harmonic
divergence cancels against the index shift and the finite head.  A direct
comparison with the already proved product quotient gives an explicit
\(C_{R,\delta}/N\) tail.

Taking the positive square root of \cref{eq:modsquare} gives the target-page
formula
\begin{align}
 |G(x+\ii y)|
 ={}&|G(x)|\exp\left(\frac{1+\gamma}{2}y^2\right)
 \sqrt{1+\frac{y^2}{x^2}}\notag\\
 &\times\sqrt{\prod_{k=1}^{\infty}
 \left(1+\frac{y^2}{(x+k)^2}\right)^{k+1}
 \exp\!\left(-\frac{y^2}{k}\right)}.
 \label{eq:modulus}
\end{align}
The squared identity is the primary iPRA theorem; the square-root corollary uses
nonnegativity and the positive square-root name.

At \(x=0\), extract the simple zero of \(G\) at the origin before taking the
limit.  One obtains
\begin{equation}\label{eq:pure-imag-square}
 |G(\ii y)|^2
 =y^2\exp((1+\gamma)y^2)
 \prod_{k=1}^{\infty}
 \left(1+\frac{y^2}{k^2}\right)^{k+1}
 \exp\!\left(-\frac{y^2}{k}\right),
\end{equation}
and therefore
\begin{equation}\label{eq:pure-imag}
 |G(\ii y)|
 =|y|\exp\left(\frac{1+\gamma}{2}y^2\right)
 \sqrt{\prod_{k=1}^{\infty}
 \left(1+\frac{y^2}{k^2}\right)^{k+1}
 \exp\!\left(-\frac{y^2}{k}\right)}.
\end{equation}
The absolute value on \(y\) is forced by both symmetry and nonnegativity.  This
is the corrected version of the source statement.

\subsection{Formalization pattern}

Reflection identities combine four reusable components: the derivative formula,
Gamma reflection, path integration, and a normalization point.  Clausen and
modulus formulas add only finite Fourier/product rearrangements.  A compact Lean
plan is therefore:

\begin{enumerate}
\item prove a generic ``zero derivative along a certified path'' telescope;
\item instantiate it with the cross-multiplied derivative of \(\Phi\);
\item define \(\Cl\) by its sine series and prove the integral form with a
finite Fourier sum plus the \(1/N\) tail;
\item derive \cref{eq:modsquare} directly at the level of finite partial
products, then pass to the common tail modulus;
\item expose unsquared formulas only through a positive-square-root contract.
\end{enumerate}

All induction is over finite lists, cutoffs or mesh cells and is therefore
quantifier-free after the fixed function symbols are introduced.
